\documentclass{beamer}
\usepackage[T1]{fontenc}
\usepackage{lmodern}
\usepackage{microtype}
% Show -- << >> as typed in code, not as a dash and guillemets.
\DisableLigatures[-,<,>]{encoding = T1, family = tt*}
\usepackage{graphicx}
\usetheme{Madrid}
\setbeamertemplate{navigation symbols}{}
%
% Example programs: doc/intro/examples. "make test" assembles and runs them.
% Every encoding on these slides was checked against niceasm and by hand.
%
\newcommand{\code}[1]{\texttt{#1}}
%
\title[NICE]{NICE -- a nice 32 bit architecture}
\subtitle{An introduction for beginners}
\author[nice\_cpu]{The nice\_cpu project}
\institute[]{The NICE architecture is by Bernd Ulmann}
\date{October 2026}
\begin{document}
 \begin{frame}
  \titlepage
 \end{frame}
%
   \begin{frame}
    \frametitle{Goal}
    Why another 32 bit processor? Why not just use a commodity one?
    \begin{itemize}
     \item First of all, there is nothing like understanding a CPU
      down to the last bit -- and then building it yourself.
     \item NICE (``NICE Is Charmingly Elegant'') was designed by
      Bernd Ulmann and described in two articles, in 1997 \cite{nice}
      and 1999 \cite{nice99}, and on a web page with an emulator in 2004
      \cite{nice_html}.
     \item In 2006, its 16 bit descendant QNICE (``Quite NICE'') followed
      \cite{qnice}. QNICE runs on FPGAs today.
     \item This project gives NICE a complete specification, an assembler,
      an emulator, a runtime library, a WebAssembly translator, and a
      first FPGA core -- and next, a faster one.
     \item This introduction explains NICE step by step. It assumes you
      know what a register and a memory address are, and nothing more.
    \end{itemize}
   \end{frame}
%
   \begin{frame}
    \frametitle{Basics}
    \begin{itemize}
     \item 32 bit data and 32 bit addresses
     \item Memory is \textbf{word addressed}: every address holds one
      32 bit word. There are no bytes.
     \item 16 registers of 32 bit, \code{R0}\dots\code{R15}
     \item One instruction word, followed by zero to three constant words
     \item Three operand instructions: \code{ADD R1, R2, R3} means
      $\code{R1} := \code{R2} + \code{R3}$
     \item Every operand can use every one of 8 addressing modes, so
      memory to memory operations are normal
     \item \textbf{Every} instruction is conditional
     \item Small instruction set: 16 ALU operations, 16 shifts, and
      3 control instructions
     \item Flags change only when an instruction asks for it
    \end{itemize}
   \end{frame}
%
   \begin{frame}
    \frametitle{Registers}
    \begin{itemize}
     \item A program sees 16 registers:
      \begin{center}
       \begin{tabular}{|c|c|c|c|c|c|c|}
        \hline
         \code{R0}&\code{R1}&\dots&\code{R12}&\code{R13}&\code{R14}&\code{R15}\\
        \hline
       \end{tabular}
      \end{center}
     \item Most are general purpose, but four are special:
      \begin{description}
       \item [\code{R0}:] Always reads as 0. Writes to it are thrown
        away. This is handy for constants and comparisons.
       \item [\code{R13}:] The stack pointer (\code{SP}), by convention
        only. The hardware treats it like any other register.
       \item [\code{R14}:] The status register (\code{SR}).
       \item [\code{R15}:] The program counter (\code{PC}). Writing to it
        is a jump.
      \end{description}
     \item Unlike QNICE, NICE has no register banks. A subroutine saves
      the registers it uses on the stack.
    \end{itemize}
   \end{frame}
%
   \begin{frame}
    \frametitle{The status register \code{R14}}
    The upper 8 bits hold the flags. The lower 24 bits (\code{VBASE})
    say where the interrupt vector table is:
    \begin{center}
     \begin{tabular}{|c|c|c|c|c|c|c|c||c|}
      \hline
       31&30&29&28&27&26&25&24&23\dots 0\\
      \hline
       \code{I}&\code{E}&\code{V}&\code{N}&\code{Z}&\code{C}&\code{X}&\code{1}&\code{VBASE}\\
      \hline
     \end{tabular}
    \end{center}
    \begin{description}
     \item [\code{1}:] Always 1. This makes ``always'' a condition.
     \item [\code{X}:] 1 if the last result was \code{0xFFFFFFFF}
     \item [\code{C}:] Carry (or borrow) flag
     \item [\code{Z}:] 1 if the last result was \code{0x00000000}
     \item [\code{N}:] 1 if the last result was negative (bit 31 set)
     \item [\code{V}:] 1 if the last operation caused a signed overflow
     \item [\code{E}:] If 1, maskable interrupts are enabled
     \item [\code{I}:] Set to 1 when an interrupt is taken
    \end{description}
   \end{frame}
%
   \begin{frame}
    \frametitle{Input/Output}
    \begin{itemize}
     \item All input and output is memory mapped, so there are no special
      I/O instructions. A device is read and written like any memory word.
     \item The upper 256 words of memory, \code{0xFFFFFF00} to
      \code{0xFFFFFFFF}, are reserved for I/O devices.
     \item The emulator has a console at \code{0xFFFFFF00}. Each word
      written there prints one character, from bits 7..0:
      \begin{center}
       \code{MOVE @0xFFFFFF00, 0x41} \quad prints \code{A}
      \end{center}
     \item The FPGA implementation will have a UART at the same address.
     \item Reading a memory address where nothing is connected returns 0.
    \end{itemize}
   \end{frame}
%
   \begin{frame}
    \frametitle{Instruction set basics}
    \begin{itemize}
     \item Almost every instruction has three operands: one destination
      and two sources, \textbf{destination first}:
      \begin{center}
       \code{SUB R1, R2, R3} \quad means \quad
       $\code{R1} := \code{R2} - \code{R3}$
      \end{center}
      QNICE has it the other way round (\code{SUB R3, R1}), so take care.
     \item The fields are called \code{dest}, \code{s0}, and \code{s1}.
     \item Operations that need only one source, such as \code{MOVE} and
      \code{NOT}, take \code{dest} and \code{s0}.
     \item Two modifiers can follow the mnemonic:
      \begin{description}
       \item [\code{[M]}:] Modify the flags. Without it, the flags stay as
        they are. \code{SUB[M] R1, R2, R3} sets \code{Z} if
        \code{R2} = \code{R3}.
       \item [\code{[C]}:] Use the carry flag as carry in (for
        \code{ADD}, \code{SUB}, \code{MDBL}, and \code{DBL}).
      \end{description}
     \item A condition can come first: \code{?Z ADD R1, R1, 1} adds 1
      only if \code{Z} is set.
    \end{itemize}
   \end{frame}
%
   \begin{frame}
    \frametitle{Instruction format}
    \begin{itemize}
     \item ALU and shift instructions share one 32 bit format:
      {\scriptsize
      \begin{center}
       \begin{tabular}{|c|c||c|c|c||c||c|c|c|}
        \hline
         31&30..28&27&26&25&24..21&20..14&13..7&6..0\\
        \hline
         \code{NEG}&\code{SEL}&\code{0}&\code{C}&\code{M}&\code{opcode}&\code{dest}&\code{s0}&\code{s1}\\
         \code{NEG}&\code{SEL}&\code{1}&\code{0}&\code{M}&\code{opcode}&\code{dest}&\code{s0}&\code{s1}\\
        \hline
       \end{tabular}
      \end{center}
      }
      \code{NEG} and \code{SEL} are the condition. Bits 27..26 are the
      class: \code{0x} ALU, \code{10} shift, \code{11} control.
     \item Each operand field is 7 bits, a mode and a register:
      \begin{center}
       \begin{tabular}{|c|c|}
        \hline
         3 bit&4 bit\\
         \code{mode}&\code{register}\\
        \hline
       \end{tabular}
      \end{center}
     \item The control instructions use this format:
      {\scriptsize
      \begin{center}
       \begin{tabular}{|c|c||c|c||c|c|c|}
        \hline
         31&30..28&27&26&25..21&20..8&7..0\\
        \hline
         \code{NEG}&\code{SEL}&\code{1}&\code{1}&\code{command}&\code{0}&\code{vector}\\
        \hline
       \end{tabular}
      \end{center}
      }
    \end{itemize}
   \end{frame}
%
   \begin{frame}
    \frametitle{Conditional execution}
    \begin{itemize}
     \item \code{SEL} picks one of the upper 8 bits of \code{SR}.
      The instruction runs only if that bit is 1, or with \code{NEG} = 1,
      only if it is 0:
      {\small
      \begin{center}
       \begin{tabular}{|c|c|l|l|}
        \hline
         \code{SEL}&Flag&Condition&With \code{NEG} = 1\\
        \hline
         \code{000}&\code{1}&\code{?1} or nothing (always)&\code{?!1} (never)\\
         \code{001}&\code{X}&\code{?X}&\code{?!X}\\
         \code{010}&\code{C}&\code{?C}&\code{?!C}\\
         \code{011}&\code{Z}&\code{?Z}&\code{?!Z}\\
         \code{100}&\code{N}&\code{?N}&\code{?!N}\\
         \code{101}&\code{V}&\code{?V}&\code{?!V}\\
         \code{110}&\code{E}&\code{?E}&\code{?!E}\\
         \code{111}&\code{I}&\code{?I}&\code{?!I}\\
        \hline
       \end{tabular}
      \end{center}
      }
     \item Because bit 24 of \code{SR} is always 1, \code{SEL} = 0 means
      ``always''.
     \item A skipped instruction does nothing at all, but the processor
      still steps over its constant words.
    \end{itemize}
   \end{frame}
%
   \begin{frame}
    \frametitle{List of instructions: ALU}
    {\small
    \begin{center}
     \begin{tabular}{|c|ll|l|}
      \hline
       Opc&Instr&Operands&Effect\\
      \hline
       \code{0}&\code{MOVE}&\code{dest, s0}&\code{dest := s0}\\
       \code{1}&\code{SUB}&\code{dest, s0, s1}&\code{dest := s0 - s1 (- C)}\\
       \code{2}&\code{MDBL}&\code{dest, s0, s1}&\code{dest := s0 + (s0 AND s1) (+ C)}\\
       \code{3}&\code{ADD}&\code{dest, s0, s1}&\code{dest := s0 + s1 (+ C)}\\
       \code{4}&\code{DBL}&\code{dest, s0}&\code{dest := s0 + s0 (+ C)}\\
       \code{5}&\code{DEC}&\code{dest, s0}&\code{dest := s0 - 1}\\
       \code{6}&\code{NOT}&\code{dest, s0}&\code{dest := NOT s0}\\
       \code{7}&\code{NOR}&\code{dest, s0, s1}&\code{dest := NOT (s0 OR s1)}\\
       \code{8}&\code{IAND}&\code{dest, s0, s1}&\code{dest := (NOT s0) AND s1}\\
       \code{9}&\code{NAND}&\code{dest, s0, s1}&\code{dest := NOT (s0 AND s1)}\\
       \code{A}&\code{XOR}&\code{dest, s0, s1}&\code{dest := s0 XOR s1}\\
       \code{B}&\code{IOR}&\code{dest, s0, s1}&\code{dest := (NOT s0) OR s1}\\
       \code{C}&\code{XNOR}&\code{dest, s0, s1}&\code{dest := NOT (s0 XOR s1)}\\
       \code{D}&\code{AND}&\code{dest, s0, s1}&\code{dest := s0 AND s1}\\
       \code{E}&\code{ONE}&\code{dest}&\code{dest := 1}\\
       \code{F}&\code{OR}&\code{dest, s0, s1}&\code{dest := s0 OR s1}\\
      \hline
     \end{tabular}
    \end{center}
    }
    The carry in, \code{(+ C)} or \code{(- C)}, is used only with
    \code{[C]}.
   \end{frame}
%
   \begin{frame}
    \frametitle{Some unusual operations}
    \begin{itemize}
     \item The operations come from the 74181, a classic ALU chip. That
      explains a few odd ones:
      \begin{description}
       \item [\code{MDBL}:] ``Masked double'': doubles \code{s0} where
        \code{s1} has ones. With \code{s1} = \code{0xFFFFFFFF} it is
        $2 \cdot \code{s0}$.
       \item [\code{IAND}, \code{IOR}:] AND and OR with \code{s0} inverted.
        \code{IAND R1, R2, R1} clears in \code{R1} the bits that are set
        in \code{R2}.
       \item [\code{ONE}:] Sets \code{dest} to 1.
      \end{description}
     \item Some things are missing on purpose:
      \begin{itemize}
       \item There is no \code{CMP}. Subtract into \code{R0} instead,
        which keeps only the flags (see ``Comparisons'').
       \item There is no \code{JMP}. Write to \code{PC} instead.
       \item There is no multiply or divide. The runtime library has
        routines for that, such as \code{\_\_mul}.
      \end{itemize}
    \end{itemize}
   \end{frame}
%
   \begin{frame}
    \frametitle{List of instructions: shifts}
    Shifts take the value in \code{s0} and the count in \code{s1}. The name
    shows where the new bit comes from and where the old bit goes:
    {\small
    \begin{center}
     \begin{tabular}{|c|l|l||c|l|l|}
      \hline
       Opc&Instr&Meaning&Opc&Instr&Meaning\\
      \hline
       \code{0}&\code{0<0}, \code{SHL}&shift in 0         &\code{8}&\code{0>0}, \code{SHR}&shift in 0\\
       \code{1}&\code{0<1}&shift in 1                     &\code{9}&\code{1>0}&shift in 1\\
       \code{2}&\code{0<C}&shift in \code{C}              &\code{A}&\code{C>0}&shift in \code{C}\\
       \code{3}&\code{C<0}&in 0, out to \code{C}          &\code{B}&\code{0>C}&in 0, out to \code{C}\\
       \code{4}&\code{C<1}&in 1, out to \code{C}          &\code{C}&\code{1>C}&in 1, out to \code{C}\\
       \code{5}&\code{C<C}&rotate through \code{C}        &\code{D}&\code{C>C}&rotate through \code{C}\\
       \code{6}&\code{<<<}, \code{ROL}&rotate             &\code{E}&\code{>>>}, \code{ROR}&rotate\\
       \code{7}&\code{C<<}, \code{BPL}&rotate, copy to \code{C}&\code{F}&\code{>>C}, \code{BPR}&rotate, copy to \code{C}\\
      \hline
     \end{tabular}
    \end{center}
    }
    \begin{itemize}
     \item Left shifts are \code{A<B}: \code{B} goes into bit 0, and bit 31
      goes to \code{A} (\code{0} means it is lost).
     \item Right shifts are \code{A>B}: \code{A} goes into bit 31, and
      bit 0 goes to \code{B}.
     \item The bit goes to \code{C} only with \code{[M]}.
    \end{itemize}
   \end{frame}
%
   \begin{frame}
    \frametitle{Shift examples}
    \begin{itemize}
     \item Multiply \code{R1} by 16:
      \begin{center}
       \code{SHL R1, R1, 4}
      \end{center}
     \item Unsigned divide \code{R1} by 2, and keep the remainder in
      \code{C}:
      \begin{center}
       \code{0>C[M] R1, R1, 1}
      \end{center}
     \item Swap the two halves of \code{R1}:
      \begin{center}
       \code{ROL R1, R1, 16}
      \end{center}
     \item There is no arithmetic shift right, but two instructions do it.
      The first copies the sign bit into \code{C} and throws the result
      away in \code{R0}. The second shifts \code{C} in from the left:
      \begin{center}
       \begin{tabular}{l}
        \code{C<0[M] R0, R1, 1}\\
        \code{C>0~~~~R1, R1, 3}\\
       \end{tabular}
      \end{center}
     \item The count uses all 32 bits of \code{s1}. A shift (but not a
      rotation) by 32 or more shifts everything out.
    \end{itemize}
   \end{frame}
%
   \begin{frame}
    \frametitle{Control instructions}
    The control class has seven commands, chosen by bits 25..21:
    \begin{center}
     \begin{tabular}{|l|l|l|l|}
      \hline
       Cmd bits&Command&Word&Description\\
      \hline
      \hline
       \code{00000}&\code{HALT}&\code{0x0C000000}&Halt the processor\\
       \code{00001}&\code{RTI}&\code{0x0C200000}&Return from interrupt\\
       \code{00010}&\code{INT n}&\code{0x0C400000} + n&Software interrupt\\
        &&&with vector n (0--255)\\
       \code{00100}--&\code{RDSPC}, \code{RDSSR},&\code{0x0C800000}--&Read or write the\\
       \code{00111}&\code{WRSPC}, \code{WRSSR}&\code{0x0CE00000} + x&saved PC or SR\\
      \hline
     \end{tabular}
    \end{center}
    \begin{itemize}
     \item All other commands are reserved. They halt the processor.
     \item Control instructions are conditional too:
      \code{?Z HALT} halts only if \code{Z} is set.
     \item They never change the flags \code{X}, \code{C}, \code{Z},
      \code{N}, or \code{V}, except that \code{RTI} restores all of
      \code{SR}, and \code{RDSPC R14} and \code{RDSSR R14} write it.
    \end{itemize}
   \end{frame}
%
   \begin{frame}
    \frametitle{Affected condition bits}
    With \code{[M]}, an instruction (W)rites the flags, or sets them to
    (0). It (R)eads \code{C} when \code{C} is used as an input:
    {\small
    \begin{center}
     \begin{tabular}{|l|c|c|c|c|c|}
      \hline
       Instruction&\code{V}&\code{N}&\code{Z}&\code{C}&\code{X}\\
      \hline
       \code{ADD}, \code{SUB}, \code{MDBL}, \code{DBL}&W&W&W&W, R with \code{[C]}&W\\
       \code{DEC}&W&W&W&W&W\\
       \code{MOVE}, \code{ONE}, and the logic operations&0&W&W&0&W\\
       Shifts that send a bit to \code{C}&0&W&W&W&W\\
       Other shifts&0&W&W&0&W\\
       \code{HALT}, \code{INT}&&&&&\\
       \code{RTI}&\multicolumn{5}{c|}{all restored}\\
      \hline
     \end{tabular}
    \end{center}
    }
    \begin{itemize}
     \item Without \code{[M]}, no flag changes.
     \item The shifts that shift in \code{C} read it, with or without
      \code{[M]}.
     \item \code{MOVE[M] R0, R5} is a handy test: \code{Z}, \code{N},
      and \code{X} then describe \code{R5}.
    \end{itemize}
   \end{frame}
%
   \begin{frame}
    \frametitle{Comparisons}
    There is no \code{CMP}, but \code{SUB[M]} into \code{R0} does the
    same job. \code{SUB[M] R0, a, b} computes $a - b$, throws the result
    away, and sets the flags:
    \begin{center}
     \begin{tabular}{|l|c|c|l|}
      \hline
       Condition&\code{Z}&\code{C}&Test\\
      \hline
       $a = b$&1&0&\code{?Z}\\
       $a \neq b$&0&&\code{?!Z}\\
       $a < b$, unsigned&0&1&\code{?C}\\
       $a \geq b$, unsigned&&0&\code{?!C}\\
      \hline
     \end{tabular}
    \end{center}
    \begin{itemize}
     \item For $a > b$, swap the operands: \code{SUB[M] R0, b, a}, then
      \code{?C}.
     \item Signed $a < b$ is \code{N} $\ne$ \code{V}. That takes more than
      one conditional instruction to test.
    \end{itemize}
   \end{frame}
%
   \begin{frame}
    \frametitle{Addressing modes: sources}
    Every operand can use one of eight addressing modes. For the sources
    \code{s0} and \code{s1}:
    {\small
    \begin{center}
     \begin{tabular}{|c|l|l|}
      \hline
       Mode&Notation&Description\\
      \hline
       \code{000}&\code{Rx}&The contents of \code{Rx}\\
       \code{001}&\code{--Rx}&Decrement \code{Rx}, then use it\\
       \code{010}&\code{Rx++}&Use \code{Rx}, then increment it\\
       \code{011}&\code{c[Rx]}&The number $\code{Rx} + c$\\
       \code{100}&\code{@Rx}&The memory word at address \code{Rx}\\
       \code{101}&\code{@--Rx}&Decrement \code{Rx}, then use the memory
        word there\\
       \code{110}&\code{@Rx++}&Use the memory word at \code{Rx}, then
        increment \code{Rx}\\
       \code{111}&\code{@c[Rx]}&The memory word at address $\code{Rx} + c$\\
      \hline
     \end{tabular}
    \end{center}
    }
    \begin{itemize}
     \item The \code{@} means ``in memory''. Mode bit 2 says the same.
     \item \code{c} is a \textbf{constant word} that follows the
      instruction word (modes \code{011} and \code{111}).
    \end{itemize}
   \end{frame}
%
   \begin{frame}
    \frametitle{Addressing modes: destinations}
    The destination \code{dest} uses the same eight modes:
    {\small
    \begin{center}
     \begin{tabular}{|c|l|l|}
      \hline
       Mode&Notation&Description\\
      \hline
       \code{000}&\code{Rx}&Write \code{Rx}\\
       \code{001}&\code{Rx--}&Write \code{Rx}, then decrement it\\
       \code{010}&\code{Rx++}&Write \code{Rx}, then increment it\\
       \code{011}&\code{c[Rx]}&Same as \code{@c[Rx]}\\
       \code{100}&\code{@Rx}&Write the memory word at address \code{Rx}\\
       \code{101}&\code{@--Rx}&Decrement \code{Rx}, then write the
        memory word there\\
       \code{110}&\code{@Rx++}&Write the memory word at \code{Rx}, then
        increment \code{Rx}\\
       \code{111}&\code{@c[Rx]}&Write the memory word at address
        $\code{Rx} + c$\\
      \hline
     \end{tabular}
    \end{center}
    }
    \begin{itemize}
     \item The sources are read first (\code{s0}, then \code{s1}), and
      the destination is written last.
     \item The destination is never read, so \code{ADD @R1, @R1, 1} reads
      the memory word once and writes it once.
    \end{itemize}
   \end{frame}
%
   \begin{frame}
    \frametitle{Using constant operands}
    A constant is a word that follows the instruction. The assembler
    writes it for you:
    \begin{itemize}
     \item \code{c[R0]} is just \code{c}, because \code{R0} is always 0.
      So this is how you write an \textbf{immediate} value:
      \begin{center}
       \code{MOVE R1, 0x1234}
      \end{center}
     \item \code{@c[R0]} is \code{@c}, a fixed (\textbf{absolute})
      memory address:
      \begin{center}
       \code{MOVE @0x40, R1}
      \end{center}
     \item With a register, it is an \textbf{indexed} address, for example
      the third word of a table that \code{R2} points to:
      \begin{center}
       \code{MOVE R1, @2[R2]}
      \end{center}
     \item One instruction can have up to three constants. They follow the
      instruction word in the order \code{dest}, \code{s0}, \code{s1}.
      \code{ADD @0x40, R1, 7} is three words long.
    \end{itemize}
   \end{frame}
%
   \begin{frame}
    \frametitle{Addressing mode examples}
    \begin{itemize}
     \item Copy \code{R2} to \code{R1}:
      \begin{center}
       \code{MOVE R1, R2}
      \end{center}
     \item Store \code{R2} in the memory word that \code{R1} points to:
      \begin{center}
       \code{MOVE @R1, R2}
      \end{center}
     \item Using \code{R13} as the stack pointer, push \code{R1} onto the
      stack, and pop it again:
      \begin{center}
       \code{MOVE @--R13, R1} \qquad \code{MOVE R1, @R13++}
      \end{center}
     \item Copy a word from one place in memory to another, and move both
      pointers on -- all in one instruction:
      \begin{center}
       \code{MOVE @R2++, @R1++}
      \end{center}
     \item Add 1 to a counter in memory:
      \begin{center}
       \code{ADD @COUNT, @COUNT, 1}
      \end{center}
    \end{itemize}
   \end{frame}
%
   \begin{frame}
    \frametitle{Jumps}
    \begin{itemize}
     \item \code{R15} is the program counter. A jump is just a write to
      \code{R15}. With \mbox{\code{PC .EQU R15}}:
      \begin{center}
       \begin{tabular}{ll}
        \code{MOVE PC, LOOP}&Jump to \code{LOOP}\\
        \code{?!Z MOVE PC, LOOP}&Jump if \code{Z} is not set\\
        \code{ADD PC, PC, 10}&Jump forward by 10 words\\
       \end{tabular}
      \end{center}
     \item When an instruction reads \code{R15}, it gets the address of the
      \textbf{next} instruction, after all constant words.
     \item So \code{ADD PC, PC, 10} jumps to 10 words after the end of the
      \code{ADD}.
     \item Any write to \code{R15} jumps, even with other addressing modes.
      For example, \code{MOVE PC, @R13++} pops an address off the stack
      and jumps there.
    \end{itemize}
   \end{frame}
%
   \begin{frame}
    \frametitle{Subroutine calls}
    \begin{itemize}
     \item There is no call instruction. A call is two instructions, with
      \mbox{\code{SP .EQU R13}}:
      \begin{center}
       \begin{tabular}{ll}
        \code{ADD ~@--SP, 2, PC}&Push the return address\\
        \code{MOVE PC, ROUTINE}&Jump to the subroutine\\
       \end{tabular}
      \end{center}
     \item Why \code{2}? \code{PC} reads as the address after the
      \code{ADD}, which is where the 2 word \code{MOVE} starts. The return
      address is after the \code{MOVE}.
     \item To return, pop the return address into the program counter:
      \begin{center}
       \code{MOVE PC, @SP++}
      \end{center}
     \item The stack grows down, and \code{SP} points at the last word
      pushed.
     \item By convention, arguments are passed in \code{R1}\dots\code{R4},
      and the result comes back in \code{R1}. \code{R7}\dots\code{R13} must
      be the same after the call as before.
    \end{itemize}
   \end{frame}
%
   \begin{frame}
    \frametitle{Examples}
    The following examples show how instructions are encoded:
    {\footnotesize\setlength{\tabcolsep}{4pt}
    \begin{center}
     \begin{tabular}{|c|c|c|c|c|c|c|c||l|}
      \hline
       \code{NEG}&\code{SEL}&class, \code{C}&\code{M}&opcode&\code{dest}&\code{s0}&\code{s1}&Hex\\
      \hline
       \multicolumn{9}{|l|}{\code{MOVE R1, R2}}\\
       \code{0}&\code{000}&\code{00}&\code{0}&\code{0000}&\code{0000001}&\code{0000010}&\code{0000000}&\code{0x00004100}\\
      \hline
       \multicolumn{9}{|l|}{\code{DEC[M] R2, R2}}\\
       \code{0}&\code{000}&\code{00}&\code{1}&\code{0101}&\code{0000010}&\code{0000010}&\code{0000000}&\code{0x02A08100}\\
      \hline
       \multicolumn{9}{|l|}{\code{?!Z MOVE R15, 4}}\\
       \code{1}&\code{011}&\code{00}&\code{0}&\code{0000}&\code{0001111}&\code{0110000}&\code{0000000}&\code{0xB003D800}\\
       \multicolumn{8}{|r||}{constant word for \code{s0}}&\code{0x00000004}\\
      \hline
       \multicolumn{9}{|l|}{\code{SHR R1, R1, 4}}\\
       \code{0}&\code{000}&\code{10}&\code{0}&\code{1000}&\code{0000001}&\code{0000001}&\code{0110000}&\code{0x090040B0}\\
       \multicolumn{8}{|r||}{constant word for \code{s1}}&\code{0x00000004}\\
      \hline
       \multicolumn{9}{|l|}{\code{ADD @--R13, 2, R15}}\\
       \code{0}&\code{000}&\code{00}&\code{0}&\code{0011}&\code{1011101}&\code{0110000}&\code{0001111}&\code{0x0077580F}\\
       \multicolumn{8}{|r||}{constant word for \code{s0}}&\code{0x00000002}\\
      \hline
     \end{tabular}
    \end{center}
    }
   \end{frame}
%
   \begin{frame}
    \frametitle{A complete example}
    One instruction can do a lot. This one is from the 2004 NICE web page:
    \begin{center}
     \code{?!Z ADD[M] @R3++, --R2, @0x12345678[R1]}
    \end{center}
    It is encoded as \code{0xB278C971}, followed by the constant
    \code{0x12345678}. If \code{Z} is 0, it does this:
    \begin{enumerate}
     \item Decrement \code{R2}, and use it as \code{s0}.
     \item Read the memory word at $\code{R1} + \code{0x12345678}$ as
      \code{s1}.
     \item Store $\code{s0} + \code{s1}$ at the address in \code{R3}, then
      increment \code{R3}.
     \item Update the flags \code{X}, \code{C}, \code{Z}, \code{N}, and
      \code{V}.
    \end{enumerate}
    If \code{Z} is 1, it does nothing, and the program continues after the
    constant word.
   \end{frame}
%
  \begin{frame}[fragile]
   \frametitle{Code example: {\scriptsize$\sum_{i=1}^{\mathtt{0x1000}}i$}}
   The file \code{doc/intro/examples/sum.nasm}, with the address and words
   of each instruction:
   {\scriptsize
   \begin{verbatim}
                            PC      .EQU    R15
                                    .ORG    0
00000000 00005800 00000000          MOVE    R1, 0
00000002 00009800 00001000          MOVE    R2, 0x1000
00000004 00604082           LOOP    ADD     R1, R1, R2
00000005 02a08100                   DEC[M]  R2, R2
00000006 b003d800 00000004      ?!Z MOVE    PC, LOOP
00000008 0c000000                   HALT
   \end{verbatim}
   }
   \begin{itemize}
    \item \code{R1} collects the sum, and \code{R2} counts down from
     \code{0x1000}.
    \item \code{DEC[M]} sets \code{Z} when \code{R2} reaches 0. Until
     then, \code{?!Z MOVE PC, LOOP} jumps back.
    \item At the end, \code{R1} = \code{0x800800} = 4096 $\cdot$ 4097 / 2.
   \end{itemize}
  \end{frame}
%
  \begin{frame}[fragile]
   \frametitle{Code example: copy a block of memory}
   {\footnotesize
   \begin{verbatim}
PC      .EQU    R15
        MOVE    R1, SRC         ; source pointer
        MOVE    R2, DST         ; destination pointer
        MOVE    R3, 5           ; word count
LOOP    MOVE    @R2++, @R1++    ; copy a word, advance both pointers
        DEC[M]  R3, R3
    ?!Z MOVE    PC, LOOP
        HALT

SRC     .DAT    1, 2, 3, 4, 5
DST     .BLK    5               ; 5 words of 0
   \end{verbatim}
   }
   \begin{itemize}
    \item Labels are words at the start of a line. Their value is an
     address.
    \item \code{.DAT} puts words into memory, and \code{.BLK} reserves
     words that start as 0.
    \item Everything after \code{;} is a comment.
   \end{itemize}
  \end{frame}
%
  \begin{frame}[fragile]
   \frametitle{Code example: Hello, NICE!}
   {\scriptsize
   \begin{verbatim}
SP      .EQU    R13
PC      .EQU    R15
CONSOLE .EQU    0xFFFFFF00

        MOVE    SP, 0x1000      ; set up the stack
        MOVE    R1, TEXT        ; the argument: the string's address
        ADD     @--SP, 2, PC    ; push the return address ...
        MOVE    PC, PUTS        ; ... and jump: a subroutine call
        MOVE    @CONSOLE, 10    ; a newline
        HALT

; PUTS prints the string at R1. The string ends with a 0 word.
PUTS    MOVE[M] R2, @R1++       ; next character, and set Z if it is 0
    ?Z  MOVE    PC, @SP++       ; at the end: return
        MOVE    @CONSOLE, R2    ; print the character
        MOVE    PC, PUTS        ; and loop

TEXT    .ASCII_L 'Hello, NICE!'
   \end{verbatim}
   }
   \code{.ASCII\_L} stores one character per word and adds a 0 word at
   the end. The conditional return \code{?Z MOVE PC, @SP++} needs no
   branch around it.
  \end{frame}
%
  \begin{frame}
   \frametitle{Interrupts}
   NICE has a simple interrupt scheme, the same as QNICE's:
   \begin{itemize}
    \item There are three kinds of interrupt: a maskable request
     (\emph{IRQ}), a non-maskable request (\emph{NMI}), and the software
     interrupt \code{INT n}.
    \item IRQs are accepted only while the flag \code{E} is 1. After
     reset it is 0.
    \item When an interrupt is taken, the processor saves \code{PC} and
     \code{SR} in two hidden registers, sets \code{I}, and jumps through
     the vector table.
    \item The interrupt service routine (\emph{ISR}) ends with \code{RTI},
     which restores \code{PC} and \code{SR}. So the flags are safe, but
     the ISR must save any other register it uses.
    \item Interrupts cannot be nested. A request that arrives during an
     ISR waits until the \code{RTI}.
   \end{itemize}
  \end{frame}
%
  \begin{frame}
   \frametitle{The vector table}
   \begin{itemize}
    \item The vector table has 256 words. Word $v$ holds the address of the
     ISR for vector $v$.
    \item The device that requests an interrupt supplies the vector number.
     \code{INT n} uses \code{n}.
    \item The table starts at address $\code{VBASE} \cdot 256$, where
     \code{VBASE} is bits 23..0 of \code{SR}.
    \item So to put the table at \code{0x100} and then allow IRQs, with
     \mbox{\code{SR .EQU R14}}:
     \begin{center}
      \begin{tabular}{ll}
       \code{OR SR, SR, 1}&\code{VBASE} = 1\\
       \code{OR SR, SR, 0x40000000}&\code{E} = 1\\
      \end{tabular}
     \end{center}
    \item After reset, \code{PC} = 0 and \code{VBASE} = 0, so word 0 is
     both the first instruction and vector table entry 0. Set \code{VBASE}
     before you allow interrupts.
   \end{itemize}
  \end{frame}
%
  \begin{frame}[fragile]
   \frametitle{Code example: an interrupt routine}
   {\fontsize{7}{8}\selectfont
   \begin{verbatim}
SP      .EQU    R13
SR      .EQU    R14
PC      .EQU    R15
CONSOLE .EQU    0xFFFFFF00

        .ORG    0
        MOVE    SP, 0x1000
        OR      SR, SR, 1               ; VBASE = 1: the table is at 0x100
        OR      SR, SR, 0x40000000      ; E = 1: accept IRQs
        MOVE    R1, 20
WAIT    DEC[M]  R1, R1                  ; count down from 20 ...
    ?!Z MOVE    PC, WAIT                ; ... while the IRQs come in
        MOVE    @CONSOLE, 10            ; a newline
        HALT

ISR     MOVE    @CONSOLE, 0x2A          ; print '*'
        RTI                             ; restore PC and SR, and return

        .ORG    0x105                   ; vector 5 of the table at 0x100
        .DAT    ISR
   \end{verbatim}
   }
   \code{niceemu --irq 5:5 --irq 12:5 irq.bin} raises two IRQs with vector
   5, and prints \code{**}.
  \end{frame}
%
  \begin{frame}
   \frametitle{Soft- and hardware}
   Everything is in the repository \code{github.com/MJoergen/nice\_cpu}:
   \begin{itemize}
    \item \code{doc/ISA.md}: The specification. When in doubt, it is
     the authority.
    \item \code{assembler/}: \code{niceasm}, the assembler, one Python 3
     script.
    \item \code{emulator/}: \code{niceemu}, the emulator, written in Zig.
    \item \code{runtime/}: Start-up code and routines such as multiply,
     divide, and floating point.
    \item \code{translator/}: \code{nicewasm}, which translates
     WebAssembly into NICE assembly. That way, C, Rust, and Zig programs
     can run on NICE.
    \item \code{hdl/}: The FPGA implementation in VHDL: a multi-cycle
     core and a pipelined core, both checked against the emulator.
   \end{itemize}
  \end{frame}
%
  \begin{frame}[fragile]
   \frametitle{Using the assembler and the emulator}
   Assemble the summation program and look at the result:
   {\footnotesize
   \begin{verbatim}
$ assembler/niceasm.py -o sum.bin doc/intro/examples/sum.nasm
$ cat sum.bin
00000000 00005800
00000001 00000000
00000002 00009800
00000003 00001000
...
   \end{verbatim}
   }
   Run it. The emulator prints the registers when it halts:
   {\footnotesize
   \begin{verbatim}
$ emulator/zig-out/bin/niceemu sum.bin
Halted (HALT) after 12291 steps, PC=00000009
R0  00000000 R1  00800800 R2  00000000 R3  00000000
R4  00000000 R5  00000000 R6  00000000 R7  00000000
R8  00000000 R9  00000000 R10 00000000 R11 00000000
R12 00000000 R13 00000000 R14 09000000 R15 00000009
   \end{verbatim}
   }
   \code{R14} = \code{0x09000000}: bit 24 is always set, and bit 27 is
   \code{Z}.
  \end{frame}
%
  \begin{frame}[fragile]
   \frametitle{Watching a program run}
   \code{--trace} prints one line per instruction: its address, its word,
   and what it changed. \code{-n} stops after a number of steps:
   {\footnotesize
   \begin{verbatim}
$ emulator/zig-out/bin/niceemu --trace -n 8 sum.bin
00000000 00005800
00000002 00009800 R2=00001000
00000004 00604082 R1=00001000
00000005 02a08100 R2=00000fff
00000006 b003d800 R15=00000004
00000004 00604082 R1=00001fff
00000005 02a08100 R2=00000ffe
00000006 b003d800 R15=00000004
Stopped at the step limit after 8 steps, PC=00000004
...
   \end{verbatim}
   }
   \begin{itemize}
    \item The first line changes nothing: \code{R1} was 0 already.
    \item \code{R15} is shown only when an instruction jumps.
    \item The FPGA implementation will write the same trace, so the two can
     be compared line by line.
   \end{itemize}
  \end{frame}
%
  \begin{frame}
   \frametitle{Things to remember}
   \begin{itemize}
    \item The destination comes first: \code{ADD dest, s0, s1}.
    \item The flags change only with \code{[M]}.
    \item Memory has words, not bytes. \code{.ASCII\_L} stores one
     character per word.
    \item \code{R0} is always 0, and \code{R15} reads as the address of the
     next instruction.
    \item A comparison is \code{SUB[M] R0, a, b}, a jump is a write to
     \code{PC}, and a call is \code{ADD @--SP, 2, PC} followed by
     \code{MOVE PC, label}.
    \item Every instruction can be conditional, so short \code{if}s need
     no jumps:
     \begin{center}
      \begin{tabular}{l}
       \code{SUB[M] ~R0, R1, R2}\\
       \code{?C~~MOVE R1, R2~~~~; R1 := max(R1, R2), unsigned}\\
      \end{tabular}
     \end{center}
   \end{itemize}
  \end{frame}
%
  \begin{frame}
   \frametitle{Future plans}
   \begin{itemize}
    \item As of now (OCT-2026) the following exists:
     \begin{itemize}
      \item a complete specification of the instruction set,
      \item an assembler and an emulator,
      \item a runtime library,
      \item a WebAssembly translator, so that C, Rust, and Zig programs
       can run on the emulator, and
      \item two FPGA cores, a simple multi-cycle one and a pipelined
       one, which write the same trace as the emulator for every test
       program.
     \end{itemize}
    \item Next comes a faster clock for the \textbf{pipelined core}.
     The system around it, with RAM and a UART, is built in the separate
     NiceMonMEGA65 repository.
    \item A C compiler that targets NICE directly would be great, too.
     The calling convention for one is already in \code{runtime/ABI.md}.
   \end{itemize}
  \end{frame}
%
 \begin{frame}
  \frametitle{Bibliography}
  {\small
  \begin{thebibliography}{9}
   \bibitem{nice}
    Bernd Ulmann, \emph{NICE -- an elegant and powerful 32-bit
    architecture}, Computer Architecture News 25(4), SEP-1997.
   \bibitem{nice99}
    Bernd Ulmann, \emph{Designing a NICE processor}, Microprocessors and
    Microsystems 23, 1999.
   \bibitem{nice_html}
    Bernd Ulmann, \emph{The NICE Processor Pages},
    \code{http://www.vaxman.de/projects/nice/nice.html}, 2004.
   \bibitem{qnice}
    Bernd Ulmann, \emph{QNICE -- a nice 16 bit architecture},
    \code{https://github.com/sy2002/QNICE-FPGA}.
   \bibitem{isa}
    \emph{NICE instruction set architecture}, \code{doc/ISA.md} in
    \code{https://github.com/MJoergen/nice\_cpu}.
  \end{thebibliography}
  }
 \end{frame}
\end{document}
